\section{The complete pole-corrected Hurwitz zero receiver}
\label{sec:hurwitz-pole-receiver}

This calculation uses the original shifted flow, not a deformation of its
coordinates. The programme source is
\nolinkurl{split_zero_cue_shifted_hurwitz_phase_space.tex},
Theorem \texttt{thm:universal-flow}, lines 602--630, and
Theorem \texttt{thm:zero-motion}, lines 1003--1060.
Its second occurrence is \nolinkurl{split_zero_secondary_note.tex},
lines 1439--1633. The original theta coordinate and factor are those of
Brad Rodgers and Terence Tao, \emph{The de Bruijn--Newman constant is
non-negative}, arXiv:1801.05914v5, equations \texttt{htdef} and \texttt{hoz}.
The elementary Hurwitz identities are recorded with original TeX encodings at
\href{https://dlmf.nist.gov/25.11.E7}{DLMF 25.11.7} and
\href{https://dlmf.nist.gov/25.11.E14}{DLMF 25.11.14}; the gamma product is
\href{https://dlmf.nist.gov/5.8.E2}{DLMF 5.8.2}.
These standard formulas are credited to their human literature; the
pole-corrected application below is a programme derivation.

\subsection{An entire receiver with an exact divisor correction}

For real $t\ge-1/2$ put $a=1+t$, and retain
\begin{equation}\tag{HPR1}
 s(z)=\frac12+\frac{iz}{2},\qquad
 A(s)=\frac{s(s-1)}2\pi^{-s/2}\Gamma(s/2),\qquad
 H_t^{\mathrm{Hur}}(z)=\frac18 A(s(z))\zeta(s(z),a).
\end{equation}
Here $H_t^{\mathrm{Hur}}$ denotes the completed \emph{Hurwitz} flow in this section.
Only $H_0^{\mathrm{Hur}}$ is the Rodgers--Tao time-zero function; their heat evolution
has a different generator. The full comparison is proved in HTB.
Define the fixed entire function
\begin{equation}\tag{HPR2}
 Q(z)=\prod_{m=1}^{\infty}\left(1+\frac{z^2}{(4m+1)^2}\right)
 =\frac{\Gamma(5/4)^2}
 {\Gamma(5/4+iz/4)\Gamma(5/4-iz/4)}.
\end{equation}
The product converges locally uniformly because the sum of the absolute
values of its nonconstant terms does. Its zeros are precisely
$\pm i(4m+1)$, each simple, and $Q(0)=1$.
For completeness, Euler's limit formula for $\Gamma$ gives
\[
 \frac{\Gamma(b)^2}{\Gamma(b+w)\Gamma(b-w)}
 =\lim_{N\to\infty}\prod_{k=0}^{N}
                 \left(1-\frac{w^2}{(b+k)^2}\right).
\]
The factors $N^{\pm w}$ cancel. Taking $b=5/4,w=iz/4$ proves
the second expression without an undetermined exponential factor.

The complete entire receiver is
\begin{equation}\tag{HPR3}
 \mathscr F_t(z)=Q(z)H_t^{\mathrm{Hur}}(z)H_t^{\mathrm{Hur}}(-z)
 =\frac{\Gamma(5/4)^2}{64\sqrt\pi},
       s(z)(1-s(z))\zeta(s(z),a)\zeta(1-s(z),a).
\end{equation}
Indeed $s(-z)=1-s(z)$, and
\[
 Q(z)=\frac{4\Gamma(5/4)^2}
 {s(1-s)\Gamma(s/2)\Gamma((1-s)/2)}.
\]
Multiplication by both factors $A/8$ proves every constant in (HPR3).
The only Hurwitz poles are simple poles at $s=1$ and $s=0$ in the
two respective factors, so the displayed polynomial cancels them.
Thus $\mathscr F_t$ is entire. It is even and real on the real $z$ axis;
more precisely $\mathscr F_t(\bar z)=\overline{\mathscr F_t(z)}$.
No zero has yet been removed from its divisor by this construction.

Also $H_t^{\mathrm{Hur}}(0)>0$ for this whole range of $t$. At $a=1/2$,
$\zeta(1/2,a)=(\sqrt2-1)\zeta(1/2)<0$.
The last sign follows by pairing successive terms in
$\eta(1/2)=\sum_{n\ge1}(-1)^{n-1}n^{-1/2}>0$ and using
$\eta(s)=(1-2^{1-s})\zeta(s)$.
For $a\ge1/2$,
$\partial_a\zeta(1/2,a)=-\zeta(3/2,a)/2<0$.
The factor $A(1/2)$ is negative. Consequently
\begin{equation}\tag{HPR4}
 H_t^{\mathrm{Hur}}(0)>0,\qquad \mathscr F_t(0)=H_t^{\mathrm{Hur}}(0)^2>0.
\end{equation}

Here are the growth and summability details needed for an \emph{all}-zero
argument. Euler--Maclaurin, with $N\ge1$, gives
\begin{align*}
 \zeta(s,a)={}&\frac{a^{1-s}}{s-1}+\frac{a^{-s}}2
 +\sum_{k=1}^{N}\frac{B_{2k}}{(2k)!}(s)_{2k-1}a^{1-s-2k}\\
 &-\frac{(s)_{2N}}{(2N)!}
       \int_0^{\infty}\widetilde B_{2N}(x)(a+x)^{-s-2N}\,dx.
\end{align*}
Its domain is $\Re s+2N>1$, excluding the displayed pole.
One proves it first for $\Re s>1$ by integrating by parts on each
unit interval, using $B_j'=jB_{j-1}$ and the matching periodic endpoint
values; the boundary at infinity vanishes. The integral is locally
normally convergent in the stated larger half-plane, giving the
continuation there. The periodic Bernoulli Fourier series gives
$\|\widetilde B_{2N}\|_\infty/(2N)!
\le 2\zeta(2N)/(2\pi)^{2N}\le4/(2\pi)^{2N}$.
For $|s|\le R$, take $N=\lceil R\rceil+2$.
Then $|(s)_j|\le(R+2N)^j$, and the absolute integral without
the periodic factor is
$a^{1-\Re s-2N}/(\Re s+2N-1)$.
Every power of the fixed positive $a$ is bounded by
$\exp(O_a(R+N))$. The finite sum has $N$ terms.
After multiplying by $s-1$ these estimates give
\begin{equation}\tag{HPR5}
 \max_{|s|\le R}|(s-1)\zeta(s,a)|
       \le \exp(C_a(R+2)\log(R+2))
\end{equation}
for a finite constant $C_a$, enlarged to cover bounded $R$.
Thus $\mathscr F_t$ has order at most one. Jensen's formula applied at its
nonzero value at zero gives $N_{\mathscr F_t}(R)=O_t(R\log(R+2))$.
Summing over annuli $2^j\le|z|<2^{j+1}$ proves
$\sum_{\mathscr F_t(w)=0}\operatorname{ord}_w(\mathscr F_t)|w|^{-2}<\infty$.

The poles of $H_t^{\mathrm{Hur}}$ are at $z=i(4m+1)$, and their residue in the
original $s$ coordinate is
\begin{equation}\tag{HPR6}
 \operatorname{Res}_{s=-2m} A(s)\zeta(s,a)
 =\frac{2(-1)^{m+1}\pi^m}{(m-1)!}B_{2m+1}(a).
\end{equation}
This follows from
$\operatorname{Res}_{s=-2m}\Gamma(s/2)=2(-1)^m/m!$
and $\zeta(-2m,a)=-B_{2m+1}(a)/(2m+1)$.
In the $z$ coordinate the residue is $-i/4$ times (HPR6),
because the function has the additional factor $1/8$ and $ds/dz=i/2$.
There is an actual simple pole precisely at the indices
\[
 I_t=\{m\ge1:B_{2m+1}(1+t)\ne0\}.
\]
At such an index, $H_t^{\mathrm{Hur}}(-i(4m+1))$ evaluates at $s=2m+1>1$;
all its factors are strictly positive. In particular the reflected
factor cannot cancel this pole. At $s=0$ and $s=1$ the function is
regular, with values $A\zeta=1/2+t$ and $1/2$, respectively.

Let $\mathcal Z_t$ be the positive divisor counting every zero of
$H_t^{\mathrm{Hur}}(z)$ and every zero of $H_t^{\mathrm{Hur}}(-z)$, with their actual multiplicities.
Let $\mathcal P_t=\sum_{m\in I_t}([i(4m+1)]+[-i(4m+1)])$.
Writing $\operatorname{div}$ for the signed divisor of a meromorphic
function, the exact equality of divisors is
\begin{equation}\tag{HPR7}
 \mathcal Z_t=\operatorname{div}(\mathscr F_t)-\operatorname{div}(Q)
                         +\mathcal P_t.
\end{equation}
It is the valuation identity for (HPR3), at every point. At an inactive
index, $Q$ supplies one artificial zero in $\mathscr F_t$, in addition to any
actual zero there; subtracting $\operatorname{div}(Q)$ removes exactly
one. At an active index, $\mathscr F_t$ is nonzero at both marked points and
the last two terms cancel. This also proves directly that the right
side has nonnegative multiplicities. In particular
\[
 \sum_{w}\mathcal Z_t(w)|w|^{-2}<\infty.
\]
These divisors are invariant under both $w\mapsto-w$ and conjugation.
The zero $w=0$ is absent by (HPR4).

\subsection{The complete corrected trace, with no discarded roots}

For $n\ge1$ define
\begin{equation}\tag{HPR8}
 \begin{split}
 S_{F,n}(t)&=\sum_{\mathscr F_t(w)=0}\operatorname{ord}_w(\mathscr F_t)w^{-2n},\\
 T_n(t)&=S_{F,n}(t)-2(-1)^n\sum_{m\ge1}(4m+1)^{-2n}
             +2(-1)^n\sum_{m\in I_t}(4m+1)^{-2n}\\
        &=\sum_w\mathcal Z_t(w)w^{-2n}.
 \end{split}
\end{equation}
All sums are absolutely convergent. The order bound, evenness and
Hadamard factorization give
$\mathscr F_t(z)=\mathscr F_t(0)\prod_{\{w,-w\}}(1-z^2/w^2)^{\operatorname{ord}_w(\mathscr F_t)}$.
Indeed the paired genus-one factors lose their exponential terms;
the remaining exponential is $e^{b z+c}$, and evenness forces $b=0$.
Logarithmic differentiation at zero consequently gives
\begin{equation}\tag{HPR9}
 \log \mathscr F_t(z)=\log \mathscr F_t(0)-\sum_{n\ge1}\frac{S_{F,n}(t)}{2n}z^{2n}
\end{equation}
on the actual zero-free disk at the origin. This fixes both the
multiplicity convention and the factor two.

Put $\mathcal R=\{z^{-1}p(z^{-2}):p\in\mathbb C[X]\}$.
The complete zero trace and its finite matrices are
\begin{equation}\tag{HPR10}
 \mathcal B_t(f,g)=\sum_w\mathcal Z_t(w)
                     \overline{f(\bar w)}g(w),\qquad
 K_d^{\rm Hur}(t)=(T_{i+j+1}(t))_{0\le i,j<d}.
\end{equation}
The summability in (HPR7) proves absolute convergence on all of
$\mathcal R$, not only on a selected finite set of zeros. The form
is Hermitian by conjugation of the divisor. Substituting the monomials
gives precisely the displayed matrix.

Here is a full proof of the all-zero criterion for this new receiver:
\begin{equation}\tag{HPR11}
 \text{every zero of }H_t^{\mathrm{Hur}}\text{ is real}
 \quad\Longleftrightarrow\quad
 K_d^{\rm Hur}(t)\text{ is positive semidefinite for every }d.
\end{equation}
The forward implication is the sum of the nonnegative terms
$\mathcal Z_t(w)|f(w)|^2$ at real $w$.
For the reverse implication collect $\pm$ pairs into reciprocal nodes
$\lambda=w^{-2}$ with positive integer masses $M_\lambda$.
They form a bounded discrete set away from its only possible accumulation
point zero, are closed under conjugation with equal masses, and satisfy
$V=\sum M_\lambda|\lambda|<\infty$.
For a polynomial $p$ the quadratic form is
\[
 \sum_\lambda M_\lambda\lambda
                     \overline{p(\bar\lambda)}p(\lambda).
\]
An off-real $w$ gives a negative real or a nonreal $\lambda_0$.
Choose $0<r<|\lambda_0|$ and the finite set
$S=\{\lambda:|\lambda|\ge r\}$.
Let $L_0(X)=\prod_{\mu\in S\setminus\{\lambda_0\}}
(X-\mu)/(\lambda_0-\mu)$, and
$B=\prod_{\mu\in S\setminus\{\lambda_0\}}
(r+|\mu|)/|\lambda_0-\mu|$.
If $\lambda_0<0$, this polynomial is real and
$p_N=(X/\lambda_0)^N L_0$ has its only nonzero value on $S$
at $\lambda_0$. The target contribution is $M_{\lambda_0}\lambda_0<0$;
the absolute value of the complete remaining tail is at most
$B^2V(r/|\lambda_0|)^{2N}$, which tends to zero.
If $\lambda_0$ is nonreal, choose $|c|=1$ with
$c^2=-\bar\lambda_0/|\lambda_0|$, and set
\[
 p_N=c(X/\lambda_0)^N L_0(X)
       +\bar c(X/\bar\lambda_0)^N\overline{L_0(\bar X)}.
\]
This is a real polynomial. Its two target contributions total
$-2M_{\lambda_0}|\lambda_0|$, while the full tail is bounded by
$4B^2V(r/|\lambda_0|)^{2N}$. Choosing a sufficiently large finite
$N$ makes the form strictly negative. These are bounds on every
remaining zero, not estimates for a sampled tail. Finally approximate
the finitely many real coefficients by rationals. On the bounded node
disk the uniform error tends to zero, and the quadratic error is at
most $V\delta(2\|p_N\|+\delta)$. Thus a rational finite test is
strictly negative too, proving the converse and a rational witness.

The calculation retains the full split base. For the original bounded
distributive lattice $L$, take
$G_L(\mathcal R)=\{(0,\lambda):\lambda\in L\}
\cup(\mathcal R\times\{1_L\})$, and set
\begin{equation}\tag{HPR12}
 \widehat{\mathcal B}_t((f,\lambda),(g,\mu))
     =(\mathcal B_t(f,g),\lambda\wedge\mu),\qquad
 \Phi(c,\nu)=(c,1_K\otimes\nu).
\end{equation}
The nonzero-amplitude carrier condition is satisfied because both
input amplitudes must then be nonzero, hence both labels top.
Additivity holds in each variable with the original join, since
meet distributes over join. Scalar multiplication and conjugation
have their original amplitudes and meet labels. Evaluation at every
zero is $(f,\lambda)\mapsto(f(w),\lambda)$; it is additive and
linear and retains $e=(0,1_L)$ at amplitude cancellation. Its
$\tau$ fibre contains only $\tau$.
The map $\Phi$ commutes with these operations by its two coordinates;
the original lattice retraction $\rho_L(1_K\otimes\nu)=\nu$
retains all intermediate labels. Formula (HPR8) is an equality
of arithmetic coefficients before this lift. It does not subtract
lattice labels, identify the independent support sectors, or invert $e$.

\subsection{The actual boundary at the unshifted fibre}

At $t=0$ every odd Bernoulli value $B_{2m+1}(1)$ vanishes.
The functional equation makes $H_0^{\mathrm{Hur}}$ even, so
\begin{equation}\tag{HPR13}
 I_0=\varnothing,\quad \mathscr F_0=Q (H_0^{\mathrm{Hur}})^2,\quad
 T_n(0)=2S_n(H_0^{\mathrm{Hur}}),\quad
 K_d^{\rm Hur}(0)=2K_d(0).
\end{equation}
In particular the arithmetic zero test at this fibre is exactly the
previous all-root receiver with the factor two retained. All labels
in (HPR12) are unchanged. Positive definiteness of $K_{64}(0)$
therefore proves only the corresponding 64-dimensional result here.

The other fibres have an explicitly different all-zero contribution.
Let $C_n=\sum_{m\ge1}(4m+1)^{-2n}$. Each Bernoulli polynomial is
nonzero, so the union $E$ of all its real roots $B_{2m+1}(1+t)=0$
is countable. Its complement is dense. On this complement $I_t$
contains every positive integer and $T_n(t)=S_{F,n}(t)$.
The latter is continuous near $t=0$ by (HPR3)--(HPR4) and the Taylor
coefficient formula (HPR9). Hence
\begin{equation}\tag{HPR14}
 \lim_{\substack{t\to0\\t\notin E}}T_n(t)
       =2S_n(H_0^{\mathrm{Hur}})+2(-1)^n C_n,\qquad
 T_n(0)=2S_n(H_0^{\mathrm{Hur}}).
\end{equation}
The explicitly computed jump is the disappearing zero--pole divisor.
There is no assertion that taking a parameter limit commutes with
discarding the poles in a meromorphic divisor.

The first corrected moment gives a stronger uniform statement.
Write
\begin{equation}\tag{HPR15}
 B(t)=\frac12\left.\frac{d^2}{ds^2}
               \log(A(s)\zeta(s,1+t))\right|_{s=1/2}.
\end{equation}
This logarithm is defined near $s=1/2$ by its positive value there.
Since $ds/dz=i/2$, differentiating (HPR3) and retaining both reflected
factors proves
\begin{equation}\tag{HPR16}
 B(t)=S_{F,1}(t)+2C_1,\qquad
 T_1(t)=B(t)-2\sum_{m\in I_t}(4m+1)^{-2}.
\end{equation}
Now $B_3(1+t)=(1+t)(t+1/2)t$. Thus $1\in I_t$ at every
$0<|t|\le1/1000$; all other terms in this last sum are nonnegative.
The reproducible interval certificate
\nolinkurl{check_hurwitz_pole_receiver.py} evaluates the entire
parameter interval $1+t\in[999/1000,1001/1000]$ at 256-bit precision
with power series of length four, using Fredrik Johansson's
\href{https://arblib.org/}{Arb} through Python-flint 0.9.0. It proves
$A(1/2)\zeta(1/2,1+t)>0$ and $B(t)<1/25$ on that interval.
It uses the coefficient of $x^2$ in
$\log(A(1/2+x)\zeta(1/2+x,1+t))$, which is exactly (HPR15),
not twice that number. No parameter samples or zero list enter.
The retained machine receipt specifies the library version, enclosure,
script hash and the independently retained theta moment comparison.
Consequently the computer-assisted rational bound is
\begin{equation}\tag{HPR17}
 \boxed{\mathcal B_t(z^{-1},z^{-1})=T_1(t)<-\frac1{25}
                \quad(0<|t|\le1/1000),\qquad T_1(0)>0.}
\end{equation}
The exact infinite-sum inequality behind it is simply
$T_1(t)\le B(t)-2/25$; none of the other poles can weaken it.
At zero, $T_1(0)=2M_1/M_0>0$ for the original positive theta moments.
An alternative argument requiring no new Hurwitz numerical enclosure
is already available: the earlier certified bounds
$M_0>3/50$ and $M_1<3/4000$ give $B(0)<1/40$.
Continuity and the same first-pole inequality prove strict negativity
on a punctured neighborhood, without specifying its radius.

The roots responsible for this jump can be located analytically.
At every original trivial zero $s=-2m$, the ordinary functional equation
gives
\[
 \zeta'(-2m)=\frac{(-1)^m(2m)!}{2(2\pi)^{2m}}\zeta(2m+1)\ne0,
 \quad
 B_{2m}=\frac{(-1)^{m+1}2(2m)!}{(2\pi)^{2m}}\zeta(2m).
\]
These follow respectively by differentiating the sine factor in the
functional equation and evaluating it at $1-2m$.
The analytic implicit function theorem therefore gives a unique
zero $\sigma_m(t)$ near $-2m$ for sufficiently small $t$, depending
on $m$. Differentiating the actual Hurwitz equation
$\partial_t\zeta(s,1+t)=-s\zeta(s+1,1+t)$ yields
\begin{equation}\tag{HPR18}
 \begin{gathered}
 \sigma_m(0)=-2m,\qquad
 \sigma_m'(0)=\frac{B_{2m}}{\zeta'(-2m)}
       =-4\frac{\zeta(2m)}{\zeta(2m+1)}<0,\\
 z_m(t)=i(4m+1)+8i\frac{\zeta(2m)}{\zeta(2m+1)}t+O_m(t^2).
 \end{gathered}
\end{equation}
The branch is real in $s$ for real $t$, by conjugation and uniqueness.
Its derivative is nonzero, so for small nonzero $t$ it leaves the
fixed gamma pole and becomes an actual off-real zero of $H_t^{\mathrm{Hur}}$.
At $t=0$ it cancels that pole exactly. The radii here are not claimed
uniform in $m$; (HPR17) separately handles the entire zero divisor
on its explicit common interval.

At the separated sheet $t=1$ there is also an exact endpoint sign
witness: $\zeta(-20)-1=-1$ and
$\zeta(-19)-1=168011/6600>0$. Thus $\zeta(s)-1$ has a real zero
in $(-20,-19)$. $A(s)$ is finite and nonzero on that open interval,
so it gives an actual zero of $H_1^{\mathrm{Hur}}$ on $i(39,41)$.
Criterion (HPR11) produces a finite rational negative trace test for
that zero, with its complete tail retained. This endpoint observation
does not assert that the simple test of (HPR17) remains negative all
the way to $t=1$.

Equations (HPR7), (HPR13) and (HPR16)--(HPR18) identify exactly which
part of the shifted all-zero trace returns to the original theta
pairing. The correction is an explicit divisor and an explicit
infinite series. Neither its nonzero shifted value nor the strict
negative value in (HPR17) determines the sign of the complete
unshifted Schur residual constructed in RWR18--27.

\subsection{Exact return to every original arithmetic moment}

The boundary correction itself gives the inverse arithmetic comparison.
Retain $J_t$ from HTB26--28, so $H_t^{\mathrm{Hur}}+J_t=H_0^{\mathrm{Hur}}$, and on the actual
nonvanishing disk about zero set
\[
 c_t=J_t/H_0^{\mathrm{Hur}},\qquad
 W_t=(1-c_t(z))(1-c_t(-z))=H_t^{\mathrm{Hur}}(z)H_t^{\mathrm{Hur}}(-z)/H_0^{\mathrm{Hur}}(z)^2.
\]
The Mellin kernel determining $J_t$ is the explicit
$d_t(x)=e^{-x}(1-e^{-tx})/(1-e^{-x})$ of HTB27.
Define the actual coefficients
\begin{equation}\tag{HPR19}
 \begin{gathered}
 E_n(t)=[z^{2n-1}]\frac{W_t'(z)}{W_t(z)},\qquad
 P_n(t)=2(-1)^n\sum_{m\in I_t}(4m+1)^{-2n},\\
 2S_n(H_0^{\mathrm{Hur}})=T_n(t)+E_n(t)-P_n(t).
 \end{gathered}
\end{equation}
To prove the last identity, differentiate
$\mathscr F_t=Q(H_0^{\mathrm{Hur}})^2W_t$, use (HPR9) for the first two entire even functions,
and compare the coefficient of $z^{2n-1}$. The result is
$S_{F,n}=S_{Q,n}+2S_n(H_0^{\mathrm{Hur}})-E_n$. Substitution in (HPR8) proves
(HPR19) for every $n$, with no limiting interchange at $t=0$.
All coefficients exist because $W_t(0)>0$.

For a finite degree $d$ let $E_d=(E_{i+j+1})$ and
$P_d=(P_{i+j+1})$, and put
$v_m=(1,-(4m+1)^{-2},\ldots,(-(4m+1)^{-2})^{d-1})^T$.
Then
\begin{equation}\tag{HPR20}
 \begin{gathered}
 2K_d(0)=K_d^{\rm Hur}(t)+E_d(t)-P_d(t),\\
 P_d(t)=-2\sum_{m\in I_t}(4m+1)^{-2}v_mv_m^T\preceq0.
 \end{gathered}
\end{equation}
The matrix series converges absolutely entrywise. Its quadratic value
on a complex vector $b$ is
$-2\sum_{m\in I_t}(4m+1)^{-2}|v_m^Tb|^2$, proving the sign.
The boundary matrix $E_d$ is retained in full and is not assigned a
sign. In particular the negative test in (HPR17) returns to the
positive original value by the exact correction $E_1-P_1$; it is
not transported by an unproved positivity-preserving map.

This also returns to the \emph{entire} earlier residual, not only
its first 64 coordinates. Put $D_n(t)=T_n(t)+E_n(t)-P_n(t)$,
$A=K_{64}(0)$ and, for every $i\ge64$,
$v_i=(D_{i+a+1}(t)/2)_{0\le a<64}$. The exact residual is
\begin{equation}\tag{HPR21}
 \mathscr S_{ij}=\frac12 D_{i+j+1}(t)-v_i^T A^{-1}v_j,
 \qquad i,j\ge64.
\end{equation}
It is independent of the chosen shift: (HPR19) identifies every entry
and every $v_i$ with its original value. Its domain remains the direct
sum of finitely supported sequences, with no upper cutoff in $i$.
The positivity of $A$ was certified in RWR15--17; the coordinate map
and inverse yielding this Schur form are proved in RWR18--27.
Their complete source is retained with this reader. Formula (HPR21)
is their exact entrywise identification through the present flow.
It does not replace the residual by the known negative pole matrix.
Lifting the equality through (HPR12) keeps a common original support
label and its tensor image on every side, including when amplitudes
cancel. Thus the boundary, pole divisor, original arithmetic trace
and complete residual have explicit connecting maps.
